\documentclass[11pt]{article}

\usepackage[a4paper,margin=1in]{geometry}
\usepackage{amsmath,amssymb,amsthm,mathtools}
\usepackage{booktabs,longtable,array}
\usepackage{enumitem}
\usepackage{microtype}
\usepackage{xcolor}
\usepackage[colorlinks=true,linkcolor=blue!55!black,citecolor=red!55!black,urlcolor=blue!55!black]{hyperref}

\newtheorem{theorem}{Theorem}[section]
\newtheorem{proposition}[theorem]{Proposition}
\newtheorem{definition}[theorem]{Definition}
\newtheorem{example}[theorem]{Example}
\newtheorem{remark}[theorem]{Remark}

% Local copies of the notation used throughout the OT4ML manuscript.
\newcommand{\RR}{\mathbb R}
\newcommand{\NN}{\mathbb N}
\newcommand{\Pp}{\mathcal P}
\newcommand{\Nn}{\mathcal N}
\newcommand{\PP}{\mathbb P}
\newcommand{\EE}{\mathbb E}
\renewcommand{\d}{\,\mathrm d}
\newcommand{\al}{\alpha}
\newcommand{\be}{\beta}
\newcommand{\de}{\delta}
\newcommand{\eps}{\epsilon}
\newcommand{\Cc}{\mathcal C}
\newcommand{\cumul}[1]{\Cc_{#1}}
\DeclareMathOperator{\Wass}{\mathcal W}
\newcommand{\Couplings}{\mathcal U}
\newcommand{\eqdef}{\coloneqq}
\newcommand{\one}{\mathbf 1}
\newcommand{\AUC}{\mathrm{AUC}}
\newcommand{\ROC}{\mathrm{ROC}}

\title{Distances Between One-Dimensional Probability Laws\\
Based on Cumulative and Quantile Functions}
\author{}
\date{}

\begin{document}

\maketitle

\begin{abstract}
This note catalogues the principal distances, discrepancies, and rank
functionals built from cumulative distribution functions or quantiles on the
real line.  It distinguishes four questions that are often conflated: whether
the object is a genuine metric, whether it is finite, whether it metrizes weak
convergence, and whether it is quantitatively comparable with a Wasserstein
distance.  The main examples are Wasserstein--Mallows quantile distances,
Kolmogorov--Smirnov and Kuiper distances, the L\'evy metric, bounded-ground
Wasserstein and Ky Fan metrics, Cram\'er and Cram\'er--von Mises
discrepancies, Anderson--Darling weights, integrated quantile and stop-loss
metrics, Lorenz curves, and ROC/AUC functionals.
\end{abstract}

\tableofcontents

\section{Setup and the four different questions}

Let $\al,\be\in\Pp(\RR)$.  Their cumulative distribution functions
(CDFs) and left-continuous generalized quantiles are
\begin{equation}\label{eq:cdf-quantile}
 \cumul{\al}(x)\eqdef\al(({-\infty},x]),
 \qquad
 q_\al(r)\eqdef\cumul{\al}^{-1}(r)
 \eqdef\inf\{x\in\RR:\cumul{\al}(x)\geq r\},
 \quad 0<r<1.
\end{equation}
The CDF uses the spatial coordinate $x$, whereas the quantile uses the mass
coordinate $r$.  This is the notation of the OT4ML manuscript.  The
distinction explains most of the different scaling and
topological behavior below.

For a proposed discrepancy $d$, four logically separate properties matter.
\begin{enumerate}[label=\textnormal{(\roman*)},leftmargin=2.2em]
\item \emph{Domain and finiteness:} identify the largest natural class on
      which $d(\al,\be)<+\infty$.
\item \emph{Metric property:} $d(\al,\be)=0$ only when
      $\al=\be$, and the triangle inequality holds.
\item \emph{Weak topology:} $d(\al_n,\al)\to0$ if and only if
      $\al_n\rightharpoonup\al$, where $\rightharpoonup$ denotes convergence in
      law, or weak (narrow) convergence against bounded continuous test
      functions; this is also often called weak-star convergence of
      probability measures.
\item \emph{Quantitative comparison:} inequalities of the form
      $d\leq C\Wass_p^a$ or $\Wass_p\leq C d^b$ hold on a specified class of laws.
\end{enumerate}
Two distances can generate the same topology without being Lipschitz
equivalent.  Conversely, a popular test statistic need not be a metric at all.
General references on probability metrics and empirical-distribution
statistics are \cite{GibbsSu2002,ShorackWellner1986,Billingsley1999}.

The basic convergence criteria are
\begin{equation}\label{eq:weak-cdf-quantile}
\begin{aligned}
 \al_n\rightharpoonup\al
 &\quad\Longleftrightarrow\quad
 \cumul{\al_n}(x)\to \cumul{\al}(x)
 \text{ at every continuity point of }\cumul{\al} \\
 &\quad\Longleftrightarrow\quad
 q_{\al_n}(r)\to q_\al(r)
 \text{ for a.e. }r\in(0,1).
\end{aligned}
\end{equation}
The last equivalence may equivalently be stated at every continuity point of
$q_\al$.  A monotone function has only countably many discontinuities.

\section{Quantile distances: Wasserstein and its variants}

\subsection{The Wasserstein--Mallows family}

The canonical $L^p$ distance between quantiles is exactly the
one-dimensional Wasserstein distance \cite{Vallender1974,Villani2009}.

\begin{proposition}[Quantile representation of Wasserstein distance]
\label{prop:quantile-wasserstein}
For $1\leq p<\infty$ and $\al,\be\in\Pp_p(\RR)$,
\begin{equation}\label{eq:wp-quantile}
 \Wass_p(\al,\be)^p
 =\int_0^1 |q_\al(r)-q_\be(r)|^p\d r.
\end{equation}
Moreover,
\begin{equation}\label{eq:winfty-quantile}
 \Wass_\infty(\al,\be)
 =\mathop{\mathrm{ess\,sup}}_{0<r<1}
   |q_\al(r)-q_\be(r)|.
\end{equation}
Consequently, for $p<\infty$,
\begin{equation}\label{eq:wp-topology}
 \Wass_p(\al_n,\al)\to0
 \quad\Longleftrightarrow\quad
 \al_n\rightharpoonup\al
 \text{ and }
 \int |x|^p\d\al_n(x)\to\int |x|^p\d\al(x).
\end{equation}
On a common compact interval, $\Wass_p$ therefore metrizes weak convergence.
\end{proposition}

Thus $\Wass_p$ does not metrize weak convergence on all of $\Pp(\RR)$: it also
controls a moment.  For example,
\begin{equation}\label{eq:outlier-example}
 \be_{\eps,R}\eqdef (1-\eps)\de_0+\eps\de_R
 \rightharpoonup\de_0\quad\text{as }\eps\to0,
 \qquad
 \Wass_p(\be_{\eps,R},\de_0)=\eps^{1/p}R,
\end{equation}
which need not tend to zero when $R\to\infty$.  The $p=\infty$ topology is
stronger still.  Even on $[0,1]$,
$\be_n=(1-1/n)\de_0+n^{-1}\de_1\rightharpoonup\de_0$ while
$\Wass_\infty(\be_n,\de_0)=1$.

\subsection{A bounded quantile metric for weak convergence}

Formula \eqref{eq:weak-cdf-quantile} gives a simple quantile metric that does
metrize weak convergence without a moment condition.

\begin{proposition}[Bounded quantile metric]
\label{prop:bounded-quantile}
The formula
\begin{equation}\label{eq:bounded-quantile}
 d_{\mathrm{q,bd}}(\al,\be)
 \eqdef\int_0^1 \min\{1,|q_\al(r)-q_\be(r)|\}\d r
\end{equation}
defines a metric on $\Pp(\RR)$ and
\begin{equation*}
 d_{\mathrm{q,bd}}(\al_n,\al)\to0
 \quad\Longleftrightarrow\quad
 \al_n\rightharpoonup\al.
\end{equation*}
\end{proposition}

\begin{proof}
The integrand is the bounded metric $\min\{1,|x-y|\}$, so the triangle
inequality holds pointwise.  Equality to zero identifies the two quantile
functions almost everywhere and hence the two laws.  Weak convergence gives
almost-everywhere convergence of quantiles by
\eqref{eq:weak-cdf-quantile}; dominated convergence then proves one
implication.  Conversely, convergence in \eqref{eq:bounded-quantile} is
convergence in measure of the quantile functions.  Every subsequence has an
almost-everywhere convergent subsubsequence, which yields weak convergence by
\eqref{eq:weak-cdf-quantile}.
\end{proof}

This is best viewed as a canonical metric for convergence in measure of
quantile functions.  Unlike \eqref{eq:wp-quantile}, it is not being asserted
to be an optimal-transport cost for the truncated ground cost: a concave or
truncated ground cost need not be minimized by the monotone coupling.

The less frequently used Ky Fan metric gives an equivalent quantile
description of the same convergence.  Applied to the canonical quantile
coupling, it reads~\cite{Billingsley1999}
\begin{equation}\label{eq:quantile-ky-fan}
 d_{\mathrm{q,KF}}(\al,\be)
 \eqdef\inf\left\{\eta>0:
 \mathrm{Leb}\bigl\{r\in(0,1):|q_\al(r)-q_\be(r)|>\eta\bigr\}
 \leq\eta\right\}.
\end{equation}
It is a metric because the exceptional sets in two successive comparisons
can be united.  It metrizes weak convergence and satisfies
\begin{equation}\label{eq:bounded-quantile-ky-fan}
 d_{\mathrm{q,bd}}\leq2d_{\mathrm{q,KF}},
 \qquad
 d_{\mathrm{q,KF}}\leq d_{\mathrm{q,bd}}^{1/2}.
\end{equation}
The first estimate splits the integral according to whether the quantile
gap exceeds $d_{\mathrm{q,KF}}$; the second is Markov's inequality.  The
numerical scale matters in \eqref{eq:quantile-ky-fan}, because its threshold
is used both as a spatial tolerance and as a probability tolerance.

\subsection{Bounded-ground Wasserstein and Fortet--Mourier metrics}

There is a related genuine transport metric that must not be confused with
the fixed quantile coupling above.  Equip $\RR$ with the bounded metric
$d_{\mathrm b}(x,y)\eqdef\min\{1,|x-y|\}$ and define
\begin{equation}\label{eq:bounded-ground-wasserstein}
 \Wass_{1,\mathrm b}(\al,\be)
 \eqdef\inf_{\pi\in\Couplings(\al,\be)}
 \int_{\RR^2}d_{\mathrm b}(x,y)\d\pi(x,y).
\end{equation}
Kantorovich--Rubinstein duality identifies this, up to harmless normalization
conventions, with the bounded-Lipschitz (or Dudley) and Fortet--Mourier
metrics~\cite{FortetMourier1953,Dudley1968,GibbsSu2002}.  It metrizes weak convergence on
$\Pp(\RR)$ and
\begin{equation}\label{eq:bounded-ground-comparisons}
 \Wass_{1,\mathrm b}(\al,\be)
 \leq d_{\mathrm{q,bd}}(\al,\be),
 \qquad
 \Wass_{1,\mathrm b}(\al,\be)
 \leq\min\{1,\Wass_1(\al,\be)\}.
\end{equation}
The first bound simply tests \eqref{eq:bounded-ground-wasserstein} with the
quantile coupling.  Equality need not hold because the truncated ground cost
is not convex in the displacement.

For comparison, the following coupling form of the Prokhorov metric is due to
Strassen~\cite{Strassen1965,Billingsley1999}:
\begin{equation}\label{eq:prokhorov-coupling}
 d_{\mathrm P}(\al,\be)
 \eqdef\inf\left\{\eta>0:\ \exists\pi\in\Couplings(\al,\be),
 \pi\{|x-y|>\eta\}\leq\eta\right\}.
\end{equation}
It also metrizes weak convergence, and direct coupling arguments give
\begin{equation}\label{eq:prokhorov-bounded-ground}
 \Wass_{1,\mathrm b}\leq2d_{\mathrm P},
 \qquad
 d_{\mathrm P}\leq\Wass_{1,\mathrm b}^{1/2},
 \qquad
 d_{\mathrm P}\leq d_{\mathrm{q,KF}}.
\end{equation}
The last inequality reflects the fact that Prokhorov optimizes over all
couplings, whereas \eqref{eq:quantile-ky-fan} fixes the monotone one.

\subsection{Weighted and trimmed quantile distances}

For a measurable weight $w:(0,1)\to[0,\infty)$, a weighted quantile distance is
\begin{equation}\label{eq:weighted-quantile}
 d_{\mathrm{q},p,w}(\al,\be)
 \eqdef\left(\int_0^1 |q_\al(r)-q_\be(r)|^p w(r)\d r\right)^{1/p}.
\end{equation}
Weights large near $0$ or $1$ emphasize tails; weights that vanish there
produce robust or trimmed comparisons.  If
$0<m\leq w(r)\leq M<\infty$ almost everywhere, then
\begin{equation}\label{eq:weighted-wp-equivalence}
 m^{1/p}\Wass_p\leq d_{\mathrm{q},p,w}\leq M^{1/p}\Wass_p,
\end{equation}
so the two metrics are bi-Lipschitz equivalent.  If $w$ vanishes on a set of
positive measure, separation may fail.  An unbounded weight near an endpoint
requires additional tail integrability.  Its topology is stronger than the
$\Wass_p$ topology if $w$ also has a positive lower bound; without that lower
bound the two topologies can instead be incomparable.

A common robust version is the central trimmed distance
\begin{equation}\label{eq:trimmed-quantile}
 d_{\mathrm{q},p,\tau}^{\mathrm{trim}}(\al,\be)
 \eqdef\left(\int_\tau^{1-\tau}
 |q_\al(r)-q_\be(r)|^p\d r\right)^{1/p},
 \qquad 0<\tau<\tfrac12.
\end{equation}
It is only a pseudometric: two laws can have identical central quantiles and
different tails.  It therefore does not metrize weak convergence on
$\Pp(\RR)$.  Optimally trimmed Wasserstein distances, in which a prescribed
fraction of mass may be discarded before transport, are related but distinct
robust-transport constructions.

\subsection{Fractional and transformed quantile metrics}

Two secondary constructions are useful when ordinary moments are not the
appropriate scale.  First, for $0<s<1$,
\begin{equation}\label{eq:fractional-quantile}
 d_{\mathrm{q},s}(\al,\be)
 \eqdef\int_0^1|q_\al(r)-q_\be(r)|^s\d r
\end{equation}
is a metric on laws with finite $s$-th absolute moment: the triangle
inequality follows from $|x-z|^s\leq|x-y|^s+|y-z|^s$.  Its convergence is
equivalent to weak convergence plus convergence of the $s$-th absolute
moments.  Indeed, convergence of the metric gives quantile convergence in
measure and
$\bigl||x|^s-|y|^s\bigr|\leq|x-y|^s$; the converse follows from almost-everywhere
quantile convergence and uniform integrability.  Despite its appearance, this
is not denoted $\Wass_s$ here.  For a
concave cost $|x-y|^s$, the monotone coupling need not minimize the transport
cost, and the OT4ML convention reserves $\Wass_p$ for $p\geq1$.

Second, let $\psi:\RR\to I$ be a strictly increasing homeomorphism onto an
interval.  The transformed-quantile metric
\begin{equation}\label{eq:transformed-quantile}
 d_{p,\psi}(\al,\be)
 \eqdef\left(\int_0^1
 |\psi(q_\al(r))-\psi(q_\be(r))|^p\d r\right)^{1/p}
 =\Wass_p(\psi_\sharp\al,\psi_\sharp\be)
\end{equation}
is finite whenever the transformed $p$-moments are finite.  A bounded choice,
such as $\psi(x)=\arctan x$, metrizes weak convergence globally.  On a common
compact interval where $0<m\leq\psi'\leq M$, it is bi-Lipschitz equivalent to
$\Wass_p$, with constants $m$ and $M$.  On $\RR_+$, the choice
$\psi(x)=\log x$ instead measures multiplicative displacement and requires a
finite $p$-th logarithmic moment; it is not equivalent to ordinary $\Wass_p$
without bounds away from zero and infinity.

\section{CDF norm distances}

\subsection{The \texorpdfstring{$L^q$}{Lq}-CDF family}

For $1\leq q<\infty$, define
\begin{equation}\label{eq:cdf-lq}
 D_q(\al,\be)
 \eqdef\left(\int_\RR |\cumul{\al}(x)-\cumul{\be}(x)|^q\d x\right)^{1/q},
\end{equation}
whenever finite, and define
\begin{equation}\label{eq:ks}
 D_\infty(\al,\be)
 \eqdef\sup_{x\in\RR}|\cumul{\al}(x)-\cumul{\be}(x)|.
\end{equation}
The special cases have different classical names:
\begin{itemize}[leftmargin=2em]
\item $D_1=\Wass_1$ is the area between the two CDFs;
\item $D_2$ is the Cram\'er distance;
\item $D_\infty$ is the Kolmogorov, uniform, or Kolmogorov--Smirnov distance.
\end{itemize}
Only $q=1$ is simultaneously an $L^q$ distance between CDFs and the same-order
$L^q$ distance between quantiles.

\begin{proposition}[Topology of finite $L^q$-CDF distances]
\label{prop:cdf-lq-topology}
Let $1\leq q<\infty$.
\begin{enumerate}[label=\textnormal{(\roman*)},leftmargin=2.2em]
\item $D_q$ is a finite metric on $\Pp_1(\RR)$.
\item $D_q(\al_n,\al)\to0$ implies
      $\al_n\rightharpoonup\al$.
\item Weak convergence alone does not imply $D_q$ convergence on the line.
\item On laws supported in a common compact interval, $D_q$ metrizes weak
      convergence.
\end{enumerate}
For $q>1$, the resulting topology on $\Pp_1(\RR)$ is strictly weaker than the
$\Wass_1$ topology and strictly stronger than weak convergence in the sense of
convergent sequences.
\end{proposition}

\begin{proof}
Since $|\cumul{\al}-\cumul{\be}|\leq1$ and the one-dimensional identity gives
$\int|\cumul{\al}-\cumul{\be}|=\Wass_1(\al,\be)$, $D_q$ is finite on
$\Pp_1(\RR)$.  The norm axioms give the triangle inequality, and equality
almost everywhere of two right-continuous CDFs implies equality everywhere.

Suppose $D_q(\al_n,\al)\to0$ and let $x$ be a continuity point of
$\cumul{\al}$.  If, along a subsequence,
$\cumul{\al_n}(x)\geq \cumul{\al}(x)+\eta$, right-continuity of $\cumul{\al}$ and
monotonicity of $\cumul{\al_n}$ produce an interval to the right of $x$ on
which the CDF difference is at least $\eta/2$.  This contradicts $L^q$
convergence.  An interval to the left handles the opposite inequality.  Thus
$\cumul{\al_n}(x)\to \cumul{\al}(x)$, proving weak convergence.

For the failure of the converse, use \eqref{eq:outlier-example} with
$\eps_n=1/n$ and $R_n=n^{q+1}$.  Then
$\be_{\eps_n,R_n}\rightharpoonup\de_0$, whereas
$D_q^q=\eps_n^qR_n=n$.  On a fixed compact interval, weak convergence
gives pointwise convergence of CDFs outside a countable set; dominated
convergence gives $D_q\to0$.

Finally, for $q>1$, choosing $R_n=n$ in the outlier example gives
$D_q^q=n^{1-q}\to0$ but $\Wass_1=1$.  The previous counterexample shows that weak
convergence need not imply $D_q$ convergence.
\end{proof}

\subsection{Exact comparison with Wasserstein distances}

The CDF family admits simple sharp-scale comparisons.

\begin{proposition}[CDF norms versus Wasserstein distances]
\label{prop:cdf-wasserstein-comparison}
Let $1\leq q<\infty$ and $1\leq p<\infty$.  Whenever the right-hand side is
finite,
\begin{equation}\label{eq:cdf-from-wasserstein}
 D_q(\al,\be)^q\leq \Wass_1(\al,\be)\leq\Wass_p(\al,\be).
\end{equation}
If both laws are supported in an interval of length $L$, then
\begin{equation}\label{eq:wasserstein-from-cdf}
 \Wass_p(\al,\be)
 \leq L^{1-1/(pq)}D_q(\al,\be)^{1/p}.
\end{equation}
For $q=\infty$, compact support gives
\begin{equation}\label{eq:wasserstein-from-ks}
 \Wass_p(\al,\be)\leq L D_\infty(\al,\be)^{1/p}.
\end{equation}
\end{proposition}

\begin{proof}
The first inequality follows from $|\cumul{\al}-\cumul{\be}|^q\leq
|\cumul{\al}-\cumul{\be}|$ and the identity $D_1=\Wass_1$.  On an interval of length
$L$, H\"older's inequality and bounded displacement give
\[
 \Wass_1\leq L^{1-1/q}D_q,
 \qquad
 \Wass_p^p\leq L^{p-1}\Wass_1.
\]
Combining them proves \eqref{eq:wasserstein-from-cdf}.  The same argument with
$\Wass_1\leq L D_\infty$ proves \eqref{eq:wasserstein-from-ks}.
\end{proof}

These are H\"older, not generally Lipschitz, comparisons.  For the translated
Diracs $\al=\de_0$ and $\be=\de_h$,
\begin{equation}\label{eq:dirac-shift-signatures}
 \Wass_p(\de_0,\de_h)=|h|,
 \qquad
 D_q(\de_0,\de_h)=|h|^{1/q},
 \qquad
 D_\infty(\de_0,\de_h)=1
 \quad(h\neq0).
\end{equation}
Thus $D_q$ and $\Wass_p$ are not generally Lipschitz equivalent when $q>1$, and
Kolmogorov distance is not even continuous with respect to $\Wass_p$ at an
atomic law.

\subsection{Cram\'er distance, energy distance, and CRPS}

The unweighted $L^2$-CDF distance should be distinguished from the
distribution-weighted Cram\'er--von Mises statistic discussed below.  If
$X,X'\sim\al$ and $Y,Y'\sim\be$ are independent, then
\begin{equation}\label{eq:cramer-energy}
 2D_2(\al,\be)^2
 =2\EE|X-Y|-\EE|X-X'|-\EE|Y-Y'|.
\end{equation}
Under the metric convention in which the right-hand side is the
\emph{squared} energy distance, the energy distance is therefore
$\sqrt2D_2$; some statistical sources call the unrooted expression itself
the energy distance~\cite{SzekelyRizzo2013}.  Hence Cram\'er distance is a
genuine metric on $\Pp_1(\RR)$, but Proposition~\ref{prop:cdf-lq-topology}
shows that it does not metrize weak convergence on the whole line and is
strictly weaker than $\Wass_1$ in its tail control.

The continuous ranked probability score (CRPS) is the proper scoring rule
\begin{equation*}
 \mathrm{CRPS}(\al,y)\eqdef\int_\RR
   \big(\cumul{\al}(x)-\one_{\{y\leq x\}}\big)^2\d x.
\end{equation*}
If $Y\sim\be$, subtracting the expected score of the true CDF gives
\begin{equation}\label{eq:crps-cramer}
 \EE\,\mathrm{CRPS}(\al,Y)
 -\EE\,\mathrm{CRPS}(\be,Y)
 =D_2(\al,\be)^2.
\end{equation}
Thus CRPS is not another pairwise distance: its excess population risk is the
squared Cram\'er distance.

\section{Supremum CDF distances}

\subsection{Kolmogorov--Smirnov and Kuiper}

The Kolmogorov--Smirnov distance is $D_\infty$ from \eqref{eq:ks}.  The
one-sided statistics
\begin{equation*}
 D^+(\al,\be)\eqdef\sup_x(\cumul{\al}(x)-\cumul{\be}(x)),
 \qquad
 D^-(\al,\be)\eqdef\sup_x(\cumul{\be}(x)-\cumul{\al}(x))
\end{equation*}
are not symmetric metrics individually.  Their maximum is the KS distance.
Their sum is the Kuiper distance
\begin{equation}\label{eq:kuiper}
 d_{\mathrm Kp}(\al,\be)\eqdef D^+(\al,\be)+D^-(\al,\be),
 \qquad
 D_\infty\leq d_{\mathrm Kp}\leq2D_\infty.
\end{equation}
Kuiper's statistic is especially natural on the circle because the sum is
insensitive to the choice of the cut point \cite{Kuiper1960}.  On the line it
has exactly the same convergence properties as KS.

There is an equivalent discrepancy interpretation.  If rays $(-\infty,x]$
in the KS definition are replaced by bounded half-open intervals, then
\begin{equation}\label{eq:interval-discrepancy}
 \begin{aligned}
 d_{\mathrm{int}}(\al,\be)
 &\eqdef\sup_{a<b}|\al((a,b])-\be((a,b])| \\
 &=\sup_{a<b}|(\cumul{\al}-\cumul{\be})(b)
                 -(\cumul{\al}-\cumul{\be})(a)|
 =d_{\mathrm Kp}(\al,\be).
 \end{aligned}
\end{equation}
Indeed, $\cumul{\al}-\cumul{\be}$ vanishes at both ends of the line, so its range is
$D^++D^-$.  Thus KS is the ray, or star-discrepancy, metric and Kuiper is the
interval-discrepancy metric.  Both are bounded above by total variation,
which takes the supremum over all Borel sets.  Total variation is much
stronger and does not metrize weak convergence: it remains one between
$\de_{1/n}$ and $\de_0$.

\begin{proposition}[Kolmogorov topology and regularity comparison]
\label{prop:ks-topology}
Kolmogorov convergence implies weak convergence.  Conversely, if
$\al_n\rightharpoonup\al$ and $\cumul{\al}$ is continuous, then
$D_\infty(\al_n,\al)\to0$.  Hence KS metrizes weak convergence at every
atomless limiting law, but not on all of $\Pp(\RR)$.

If $\al$ and $\be$ have densities bounded by $M$ and $p\geq1$, then
\begin{equation}\label{eq:ks-from-wp-density}
 D_\infty(\al,\be)
 \leq (p+1)p^{-p/(p+1)}
 M^{p/(p+1)}\Wass_p(\al,\be)^{p/(p+1)}.
\end{equation}
\end{proposition}

\begin{proof}
The first statement is P\'olya's uniform-convergence theorem for distribution
functions; the counterexample at an atom is
$\de_{1/n}\rightharpoonup\de_0$ with KS distance one.

For the quantitative bound, couple $X\sim\al$ and $Y\sim\be$.  For
$h>0$, monotonicity and the density bound give
\[
 \cumul{\al}(x)-\cumul{\be}(x)
 \leq \PP(|X-Y|>h)+\cumul{\be}(x+h)-\cumul{\be}(x)
 \leq \frac{\EE|X-Y|^p}{h^p}+Mh.
\]
Interchanging the laws controls the opposite difference.  Take an optimal
coupling and minimize $\Wass_p^p h^{-p}+Mh$ over $h>0$.
\end{proof}

The density assumption is essential: the translated-Dirac example
\eqref{eq:dirac-shift-signatures} rules out every bound
$D_\infty\leq C\Wass_p^a$ with $a>0$ on arbitrary laws.  Conversely,
\eqref{eq:wasserstein-from-ks} shows that KS controls $\Wass_p$ in H\"older form
on a common compact support.

\subsection{The L\'evy metric}

The L\'evy metric softens both the horizontal and vertical comparison of
CDFs:
\begin{equation}\label{eq:levy}
 d_{\mathrm L}(\al,\be)
 \eqdef\inf\left\{\eps>0:
 \cumul{\al}(x-\eps)-\eps
 \leq \cumul{\be}(x)
 \leq \cumul{\al}(x+\eps)+\eps
 \text{ for every }x\right\}.
\end{equation}
Unlike KS, it permits a small horizontal displacement of an atom.

\begin{proposition}[L\'evy metric and weak convergence]
\label{prop:levy}
The L\'evy metric metrizes weak convergence on all of $\Pp(\RR)$.  Moreover,
\begin{equation}\label{eq:levy-comparisons}
 d_{\mathrm L}(\al,\be)\leq D_\infty(\al,\be),
 \qquad
 d_{\mathrm L}(\al,\be)
 \leq \min\{1,\Wass_p(\al,\be)^{p/(p+1)}\}.
\end{equation}
There is no converse control of $\Wass_p$ by $d_{\mathrm L}$ on an unbounded
space without a moment or tail assumption.
\end{proposition}

\begin{proof}
The metrization statement is classical \cite{Levy1925,Billingsley1999}; the
first inequality is immediate from the definitions.  For the second, an
optimal $p$-transport coupling and Markov's inequality give
\[
 \PP(|X-Y|>\eps)\leq\Wass_p^p/\eps^p.
\]
Taking $\eps=\Wass_p^{p/(p+1)}$ and using the coupling characterization of
the Prokhorov metric yields $d_{\mathrm L}\leq d_{\mathrm P}\leq\eps$.
Finally, the outlier family \eqref{eq:outlier-example} can converge weakly,
and hence in L\'evy metric, while its $p$-Wasserstein distance stays bounded
away from zero or diverges.
\end{proof}

On a common compact interval, L\'evy and every finite $\Wass_p$ generate the same
topology, because both metrize weak convergence there.  This topological
equivalence should not be mistaken for a scale-free global Lipschitz
equivalence.

\section{Weighted CDF discrepancies and goodness-of-fit statistics}

\subsection{Fixed weighted CDF metrics}

Let $\nu$ be a finite Borel measure on $\RR$.  The fixed-reference weighted
family is
\begin{equation}\label{eq:fixed-weighted-cdf}
 D_{q,\nu}(\al,\be)
 \eqdef\left(\int_\RR |\cumul{\al}-\cumul{\be}|^q\d\nu\right)^{1/q}.
\end{equation}
Replacing Lebesgue measure by a finite reference removes the tail obstruction
of Proposition~\ref{prop:cdf-lq-topology}.

\begin{proposition}[A weighted CDF metric for weak convergence]
\label{prop:weighted-cdf-weak}
Assume that $\nu$ is finite, nonatomic, and gives positive mass to every
nonempty open interval.  Then $D_{q,\nu}$ is a metric on $\Pp(\RR)$ and
metrizes weak convergence for every $1\leq q<\infty$.
\end{proposition}

\begin{proof}
Full support and right-continuity ensure separation.  If
$\al_n\rightharpoonup\al$, the CDFs converge outside the countable set of
discontinuities of $\cumul{\al}$.  This set is $\nu$-null, so dominated
convergence applies.  Conversely, the monotonicity argument from
Proposition~\ref{prop:cdf-lq-topology} works because every interval has
positive $\nu$-mass.
\end{proof}

For example, one may take $\nu$ to be a Gaussian or logistic law.  This metric
depends on the chosen coordinate system and reference measure.  If
$\d\nu=w(x)\d x$ with $\|w\|_\infty<\infty$, then
\begin{equation}\label{eq:weighted-cdf-w1}
 D_{q,\nu}^q\leq\|w\|_\infty\Wass_1.
\end{equation}
There is generally no reverse moment control: a reference with rapidly
decaying tails deliberately makes remote discrepancies inexpensive.

\subsection{Cram\'er--von Mises, Anderson--Darling, and Watson}

Classical goodness-of-fit statistics usually integrate against the reference
distribution rather than against Lebesgue measure.  Taking $\be$ as the
target, the population Cram\'er--von Mises (CvM) discrepancy is
\begin{equation}\label{eq:cvm}
 \mathrm{CvM}_{\be}(\al)^2
 \eqdef\int_\RR(\cumul{\al}(x)-\cumul{\be}(x))^2\d\be(x).
\end{equation}
If $\cumul{\be}$ is continuous, then
\begin{equation}\label{eq:cvm-pp}
 \mathrm{CvM}_{\be}(\al)^2
 =\int_0^1\big(\cumul{\al}(q_\be(r))-r\big)^2\d r.
\end{equation}
The expression is the $L^2$ deviation of the probability--probability (P--P),
or ordinal-dominance, curve from the diagonal.  It is distribution-free under
the null after the probability integral transform
\cite{ShorackWellner1986}.  This is not the Cram\'er distance $D_2$, which uses
$\d x$ and is tied to energy distance by \eqref{eq:cramer-energy}.

The Anderson--Darling discrepancy gives more weight to both tails:
\begin{equation}\label{eq:ad}
 \mathrm{AD}_{\be}(\al)^2
 \eqdef\int_\RR
 \frac{(\cumul{\al}(x)-\cumul{\be}(x))^2}
 {\cumul{\be}(x)(1-\cumul{\be}(x))}\d\be(x)
 =\int_0^1\frac{|\cumul{\al}(q_\be(r))-r|^2}{r(1-r)}\d r,
\end{equation}
when the expressions are finite \cite{AndersonDarling1952}.  Watson's
$U^2$ statistic is a centered CvM statistic used for circular data; it removes
the mean vertical displacement of the P--P curve so that the statistic is
insensitive to the arbitrary circular origin.

These objects require care when called ``distances.''  The one-sample forms
\eqref{eq:cvm}--\eqref{eq:ad} are asymmetric because $\be$ is privileged.  They
can fail to separate laws when $\be$ does not have sufficiently rich support,
and the Anderson--Darling value can be infinite.  Two-sample versions often
replace $\be$ by the pooled law $(\al+\be)/2$; these are symmetric test
statistics, but their pair-dependent integration measure does not
automatically yield a triangle inequality.  They also do not metrize weak
convergence globally: $\de_{1/n}\rightharpoonup\de_0$ is again a
counterexample for statistics that inspect the jump of the atomic reference.
At a continuous limiting CDF, dominated versions behave much more like KS and
CvM empirical-process statistics.

\section{Integrated CDFs and integrated quantiles}

\subsection{Stop-loss metric}

For $\al\in\Pp_1(\RR)$, its stop-loss, call, or integrated-survival
transform is
\begin{equation}\label{eq:stoploss-transform}
 S_\al(t)\eqdef\int_\RR(x-t)_+\d\al(x)
 =\int_t^\infty(1-\cumul{\al}(x))\d x.
\end{equation}
It is a convex function whose one-sided derivatives recover the survival
function, so it determines the law.  The stop-loss metric is
\begin{equation}\label{eq:stoploss-metric}
 d_{\mathrm{SL}}(\al,\be)
 \eqdef\sup_{t\in\RR}|S_\al(t)-S_\be(t)|.
\end{equation}
It is classical in stochastic-order and actuarial theory
\cite{MullerStoyan2002,DenuitEtAl2005}.

\begin{proposition}[Stop-loss and $\Wass_1$ topologies]
\label{prop:stoploss-w1}
The stop-loss formula defines a metric on $\Pp_1(\RR)$ and
\begin{equation}\label{eq:stoploss-upper}
 d_{\mathrm{SL}}(\al,\be)\leq\Wass_1(\al,\be).
\end{equation}
Moreover, $d_{\mathrm{SL}}(\al_n,\al)\to0$ if and only if
$\Wass_1(\al_n,\al)\to0$.  Thus the two metrics induce the same topology,
but they are not globally bi-Lipschitz equivalent.
\end{proposition}

\begin{proof}
Every function $x\mapsto(x-t)_+$ is $1$-Lipschitz, so
\eqref{eq:stoploss-upper} follows from Kantorovich--Rubinstein duality.
The transform determines its law through its one-sided derivatives, proving
separation; the other metric axioms are immediate.

Suppose $d_{\mathrm{SL}}(\al_n,\al)\to0$.  Uniform convergence of the
convex functions $S_{\al_n}$ implies convergence of their derivatives at
every differentiability point of $S_\al$, hence weak convergence of the
laws.  In addition, $S_{\al_n}(0)\to S_\al(0)$ gives convergence of the
positive first moments.  Finally,
$S_\al(t)+t\to\int x\d\al(x)$ as $t\to-\infty$, and uniform convergence
of the transforms gives convergence of the means.  Positive and negative
first moments therefore converge, which together with weak convergence is
equivalent to $\Wass_1$ convergence.  The converse follows from
\eqref{eq:stoploss-upper}.

For failure of a reverse Lipschitz bound, put equal masses on the pairs
$\{4k,4k+3\}$ for $\al_n$ and $\{4k+1,4k+2\}$ for $\be_n$,
$k=0,\ldots,n-1$.  Their CDF difference alternates in sign on each block,
giving $\Wass_1(\al_n,\be_n)=1$ but
$d_{\mathrm{SL}}(\al_n,\be_n)=1/(2n)$.
\end{proof}

The stop-loss metric is therefore a useful example of topological equivalence
without Lipschitz equivalence.  It is also a natural metric for convex-order
questions because $\al\preceq_{\mathrm{cx}}\be$ is characterized, under
equality of means, by $S_\al(t)\leq S_\be(t)$ for every $t$.

\subsection{Integrated quantiles and Lorenz curves}

The lower integrated quantile, or generalized Lorenz curve, is
\begin{equation}\label{eq:integrated-quantile}
 M_\al(r)\eqdef\int_0^r q_\al(s)\d s,
 \qquad 0\leq r\leq1.
\end{equation}
This is the notation already used for integrated quantiles in the OT4ML
manuscript.  It determines $q_\al$ by differentiation almost everywhere and
hence determines $\al$.

\begin{proposition}[Integrated-quantile metric]
\label{prop:integrated-quantile-metric}
The formula
\begin{equation}\label{eq:integrated-quantile-metric}
 d_{\mathrm{IQ}}(\al,\be)
 \eqdef\sup_{0\leq r\leq1}|M_\al(r)-M_\be(r)|
\end{equation}
defines a metric on $\Pp_1(\RR)$, and
\begin{equation}\label{eq:integrated-quantile-w1}
 d_{\mathrm{IQ}}(\al,\be)
 \leq\int_0^1|q_\al-q_\be|\d r
 =\Wass_1(\al,\be),
\end{equation}
while $d_{\mathrm{IQ}}(\al_n,\al)\to0$ if and only if
$\Wass_1(\al_n,\al)\to0$.  The two metrics are not globally bi-Lipschitz
equivalent.
\end{proposition}

\begin{proof}
Only the converse topological implication requires explanation.  Uniform
convergence of the convex functions $M_{\al_n}$ implies convergence of their
derivatives at every differentiability point of $M_\al$, hence almost-everywhere
convergence of $q_{\al_n}$ to $q_\al$ and weak convergence of the laws.  In
addition,
\begin{equation*}
 \int x_-\d\al(x)=-\min_{0\leq r\leq1}M_\al(r),
 \qquad
 \int x_+\d\al(x)=M_\al(1)-\min_{0\leq r\leq1}M_\al(r).
\end{equation*}
Uniform convergence therefore gives convergence of both first-moment tails,
which upgrades weak convergence to $\Wass_1$ convergence.  The alternating
atomic example in the proof of Proposition~\ref{prop:stoploss-w1} has
$\Wass_1=1$ but $d_{\mathrm{IQ}}=1/(2n)$, ruling out a reverse Lipschitz bound.
\end{proof}

Integrated quantiles are convex conjugates of integrated CDF or stop-loss
transforms, and are standard tools for convex and increasing-convex stochastic
orders.

For a nonnegative law with mean $m_\al>0$, the normalized Lorenz curve is
\begin{equation}\label{eq:lorenz}
 L_\al(r)\eqdef\frac{M_\al(r)}{m_\al}.
\end{equation}
Distances such as $\|L_\al-L_\be\|_1$ or
$\|L_\al-L_\be\|_\infty$ compare inequality profiles
\cite{Lorenz1905,Gastwirth1971}.  They are not metrics on probability laws:
$L_{(x\mapsto ax)_\sharp\al}=L_\al$ for every $a>0$.  They become
metrics only after quotienting by positive scale or after retaining the mean
as an additional coordinate.  Consequently, a Lorenz-curve distance alone
cannot metrize weak convergence or be equivalent to $\Wass_p$ on the original
line.

The Gini coefficient is twice the area between one Lorenz curve and the
diagonal.  It is a scalar summary of one law, not a distance between two laws.

\subsection{Higher integrated CDFs and Zolotarev metrics}

Zolotarev's ideal metrics extend the Kantorovich--Rubinstein viewpoint by
testing against functions with higher-order H\"older derivatives
\cite{Zolotarev1978}.  If $s=m+\gamma$, where $m\in\NN_0$ and
$0<\gamma\leq1$, one considers
\begin{equation}\label{eq:zolotarev}
 \zeta_s(\al,\be)
 \eqdef\sup_f\left|\int f\d(\al-\be)\right|,
 \qquad
 |f^{(m)}(x)-f^{(m)}(y)|\leq|x-y|^\gamma.
\end{equation}
For $s=1$, this is $\Wass_1$.  For $m\geq1$, finiteness requires the moments
through order $m$, which are invisible to the H\"older seminorm, to agree;
finite $s$-th absolute moments are a standard sufficient integrability
assumption.  Repeated integration by parts then represents the metric through
iterated primitives of $\cumul{\al}-\cumul{\be}$, generalizing the stop-loss
construction.  These are moment-sensitive approximation metrics, not metrics
for weak convergence on all of $\Pp(\RR)$.

\section{P--P curves, ROC curves, and AUC}

\subsection{Ordinal-dominance and ROC representations}

Assume first that $\cumul{\al_0}$ is continuous and strictly increasing.  The
ordinal-dominance, or P--P, curve of $\al_1$ relative to $\al_0$ is
\begin{equation}\label{eq:odc}
 \mathrm{ODC}_{0,1}(u)\eqdef\cumul{\al_1}(q_{\al_0}(u)),
 \qquad 0<u<1.
\end{equation}
For score laws $\al_0$ in the negative class and $\al_1$ in the positive
class, one common ROC convention is
\begin{equation}\label{eq:roc}
 \ROC_{0,1}(u)\eqdef 1-\cumul{\al_1}(q_{\al_0}(1-u)),
 \qquad 0<u<1,
\end{equation}
where $u$ is the false-positive rate.  ROC is a reflected ODC curve
\cite{Bamber1975,GneitingVogel2022}.

Several CDF discrepancies are simply norms of this curve's deviation from the
diagonal:
\begin{align}\label{eq:rank-discrepancies}
 D_\infty(\al_0,\al_1)
 &=\sup_{0<u<1}|\mathrm{ODC}_{0,1}(u)-u|,\\
 \mathrm{CvM}_{\al_0}(\al_1)^2
 &=\int_0^1|\mathrm{ODC}_{0,1}(u)-u|^2\d u,\\
 \mathrm{AD}_{\al_0}(\al_1)^2
 &=\int_0^1\frac{|\mathrm{ODC}_{0,1}(u)-u|^2}{u(1-u)}\d u.
\end{align}
Thus KS is the maximal vertical P--P deviation, CvM its average squared
deviation, and Anderson--Darling a tail-weighted version.

These rank constructions are invariant when the same strictly increasing
transformation is applied to both score distributions.  Wasserstein distances
are instead sensitive to the physical scale.  Therefore no two-sided
Lipschitz or H\"older equivalence with $\Wass_p$ can hold without fixing a scale
and imposing regularity bounds.

\subsection{Area under the ROC curve is not a probability metric}

For independent $X_0\sim\al_0$ and $X_1\sim\al_1$, the tie-corrected
area under the ROC curve is
\begin{equation}\label{eq:auc}
 \AUC(\al_0,\al_1)
 \eqdef\PP(X_1>X_0)+\tfrac12\PP(X_1=X_0).
\end{equation}
For continuous laws this is
$\PP(X_1>X_0)=\int \cumul{\al_0}(y)\d\al_1(y)$ and equals
the population Mann--Whitney--Wilcoxon ranking functional
\cite{MannWhitney1947,Bamber1975}.

Neither $\AUC$ nor $|\AUC-1/2|$ is a distance between laws.  It is oriented,
does not satisfy a triangle inequality, and fails separation.  For example,
if $\al_0$ and $\al_1$ are two different centered symmetric laws, then
$X_1-X_0$ is symmetric and $\AUC=1/2$ (when ties have zero probability).
Taking $\al_0=\Nn(0,1)$ and
$\al_1=\Nn(0,4)$ gives distinct laws with zero AUC deviation from
$1/2$.  AUC is also invariant under every common increasing reparameterization
of scores, whereas $\Wass_p$ changes under nonlinear reparameterization and
scales linearly under dilation.

For laws with finite first moments, three different ``areas'' must therefore
not be confused:
\begin{enumerate}[label=\textnormal{(\roman*)},leftmargin=2.2em]
\item $\int|\cumul{\al}-\cumul{\be}|\d x=\Wass_1(\al,\be)$ is the absolute area
      between CDFs and is a metric;
\item $\int(\cumul{\al}-\cumul{\be})\d x
      =\int x\d\be-\int x\d\al$ is signed area and only compares means;
\item area under an ROC curve is the ranking probability
      \eqref{eq:auc}, not a transport distance.
\end{enumerate}

\section{A compact comparison catalogue}

The following table summarizes the main objects.  ``Weak globally'' means on
all of $\Pp(\RR)$; moment-restricted statements are indicated separately.

\renewcommand{\arraystretch}{1.25}
\begin{longtable}{@{}>{\raggedright\arraybackslash}p{.18\textwidth}
                  >{\raggedright\arraybackslash}p{.25\textwidth}
                  >{\raggedright\arraybackslash}p{.20\textwidth}
                  >{\raggedright\arraybackslash}p{.27\textwidth}@{}}
\toprule
Object & Formula & Metric status & Convergence and relation to Wasserstein \\
\midrule
\endfirsthead
\toprule
Object & Formula & Metric status & Convergence and relation to Wasserstein \\
\midrule
\endhead
Wasserstein--Mallows $\Wass_p$, $p<\infty$
& $\|q_\al-q_\be\|_{L^p(0,1)}$
& Metric on $\Pp_p(\RR)$
& Metrizes weak convergence plus $p$-moment convergence; weak on common compact support. \\

$\Wass_\infty$
& $\|q_\al-q_\be\|_\infty$
& Extended metric
& Strictly stronger than weak convergence, even on compact support. \\

Bounded quantile metric
& $\int_0^1\min(1,|q_\al-q_\be|)$
& Metric on $\Pp(\RR)$
& Metrizes weak convergence globally. \\

Quantile Ky Fan
& $\inf\{\eta:\mathrm{Leb}(|q_\al-q_\be|>\eta)\leq\eta\}$
& Metric on $\Pp(\RR)$
& Metrizes weak convergence; quantitatively equivalent to the bounded quantile metric as in \eqref{eq:bounded-quantile-ky-fan}. \\

Bounded-ground $\Wass_{1,\mathrm b}$
& OT for $d_{\mathrm b}=\min(1,|x-y|)$
& Metric on $\Pp(\RR)$
& Fortet--Mourier type; metrizes weak convergence and is no larger than the bounded quantile metric. \\

Prokhorov
& Best coupling with $\pi\{|x-y|>\eta\}\leq\eta$
& Metric on $\Pp(\RR)$
& Metrizes weak convergence; related to $\Wass_{1,\mathrm b}$ by \eqref{eq:prokhorov-bounded-ground}. \\

Weighted quantile $d_{\mathrm{q},p,w}$
& $\big(\int|q_\al-q_\be|^p w\big)^{1/p}$
& Metric if $w>0$ a.e.; possibly extended
& Bi-Lipschitz to $\Wass_p$ if $w$ is bounded above and below; otherwise changes tail sensitivity. \\

Trimmed quantile
& $\big(\int_\tau^{1-\tau}|q_\al-q_\be|^p\big)^{1/p}$
& Pseudometric
& Ignores tails; does not metrize weak convergence. \\

Fractional quantile $d_{\mathrm{q},s}$
& $\int_0^1|q_\al-q_\be|^s$, $0<s<1$
& Metric under finite $s$-moments
& Metrizes weak convergence plus $s$-moment convergence; it is not the concave-cost OT value. \\

Transformed quantile $d_{p,\psi}$
& $\|\psi\circ q_\al-\psi\circ q_\be\|_{L^p}$
& Metric under transformed moments
& A bounded homeomorphism $\psi$ gives the weak topology; local derivative bounds give comparison with $\Wass_p$. \\

$L^q$ CDF, $q<\infty$
& $D_q=\|\cumul{\al}-\cumul{\be}\|_{L^q(\d x)}$
& Metric on $\Pp_1(\RR)$
& $D_q^q\leq\Wass_1$; implies weak convergence, but weak convergence does not imply $D_q$ convergence without tail control. \\

Cram\'er / energy
& $D_2$; squared energy${}=2D_2^2$
& Metric on $\Pp_1(\RR)$
& Strictly weaker in tails than $\Wass_1$; weak on common compact support. \\

Kolmogorov--Smirnov
& $\|\cumul{\al}-\cumul{\be}\|_\infty$
& Metric
& Implies weak; converse at continuous limiting CDFs, but not at atoms. No control by $\Wass_p$ without regularity. \\

Kuiper
& $\sup(F-G)+\sup(G-F)$
& Metric; equivalent to KS
& Same topology as KS; equals interval discrepancy and is adapted to circular data. \\

L\'evy
& Horizontal and vertical CDF tolerance
& Metric on $\Pp(\RR)$
& Metrizes weak convergence globally; $d_{\mathrm L}\lesssim\Wass_p^{p/(p+1)}$, but no reverse tail control. \\

Fixed weighted CDF
& $\|\cumul{\al}-\cumul{\be}\|_{L^q(\nu)}$
& Metric if $\nu$ has full support
& Metrizes weak convergence if $\nu$ is finite, nonatomic, and has full support. \\

CvM / Anderson--Darling
& $\int(F-G)^2\d G$; tail-weighted analogue
& Usually asymmetric discrepancies, not metrics
& Rank-based test statistics; no global weak metrization in their usual reference-dependent form. \\

Stop-loss
& $\sup_t|\EE(X-t)_+-\EE(Y-t)_+|$
& Metric on $\Pp_1(\RR)$
& Same topology as $\Wass_1$ and bounded above by $\Wass_1$, but not bi-Lipschitz equivalent. \\

Integrated quantile
& $\sup_{r\in[0,1]}|M_\al(r)-M_\be(r)|$
& Metric on $\Pp_1(\RR)$
& Same topology as $\Wass_1$ and bounded above by $\Wass_1$, but not bi-Lipschitz equivalent. \\

Zolotarev $\zeta_s$
& Smooth test functions; equivalently iterated CDF primitives
& Metric on moment-matched classes
& $\zeta_1=\Wass_1$; higher orders impose moment matching and do not metrize weak convergence globally. \\

Lorenz-curve norm
& $\|L_\al-L_\be\|$
& Pseudometric on laws; metric modulo scale
& Scale invariant, hence cannot metrize weak convergence of unnormalized laws or be equivalent to $\Wass_p$. \\

ROC/AUC
& $\PP(X_1>X_0)+\frac12\PP(X_1=X_0)$
& Ranking functional, not a metric
& Monotone-transform invariant and non-separating; no Wasserstein equivalence and no induced topology on laws. \\
\bottomrule
\end{longtable}

\section{Two counterexamples that diagnose most claims}

The following two families are useful consistency checks for any proposed
comparison.

\paragraph{Small translation.}
For $\de_0$ and $\de_h$, formula
\eqref{eq:dirac-shift-signatures} gives
$\Wass_p=|h|$, $D_q=|h|^{1/q}$, and KS${}=1$.  Thus spatial transport metrics
see the size of the displacement, finite CDF norms see both displacement and
the selected exponent, and rank suprema see only that a jump moved.

\paragraph{Rare remote outlier.}
For $\de_0$ and $\be_{\eps,R}$ from
\eqref{eq:outlier-example},
\begin{equation}\label{eq:outlier-signatures}
 \Wass_p=\eps^{1/p}R,
 \qquad
 D_q=\eps R^{1/q},
 \qquad
 D_\infty=\eps,
 \qquad
 d_{\mathrm L}\to0\quad\text{when }\eps\to0.
\end{equation}
The free choices of $\eps$ and $R$ show exactly which quantities
control tail location, tail mass, or only weak convergence.

\section{Practical reading guide}

\begin{itemize}[leftmargin=2em]
\item Use $\Wass_p$ when the geometry and the size of displacement matter and a
      $p$-moment assumption is meaningful.
\item Use L\'evy, Prokhorov, or the bounded quantile metric
      \eqref{eq:bounded-quantile} when the sole target is convergence in law.
      The bounded-ground Wasserstein metric is the OT/Fortet--Mourier version;
      the Ky Fan formula is useful when probability and spatial tolerances
      should be shown separately.
\item Use KS or Kuiper for uniform rank discrepancies, remembering that atoms
      make them strictly stronger than weak convergence.
\item Use Cram\'er/energy distance or CRPS when an integrated squared CDF
      error is desired; do not confuse it with the reference-weighted CvM
      statistic.
\item Use CvM or Anderson--Darling for goodness-of-fit testing after a
      probability-integral transform; they are not generic transport metrics.
\item Use weighted or trimmed quantile norms when tail emphasis or robustness
      is intentional, and state the weight because it determines the topology.
      Transformed quantiles are preferable when displacement is naturally
      measured after a monotone change of scale, such as a logarithm.
\item Use stop-loss or integrated quantiles for convex-order and actuarial
      questions.  Use normalized Lorenz curves only when scale invariance is
      intended.
\item Treat ROC and AUC as two-sample ranking summaries, not as distances on
      the Wasserstein space.
\end{itemize}

\bibliographystyle{plain}
\bibliography{references}

\end{document}
